% preprint/modules/discussion.tex
% Shared-Horizon Preprint — Narrative Modules
% Module: Discussion

\section{Discussion}
\label{sec:discussion}

The Shared-Horizon framework provides a mathematically rigorous and conceptually
clear method for comparing diverse AI lineages without inducing collapse between
their internal geometries. By integrating operator geometry, expressive
manifolds, reversible narrative operators, stability envelopes, transition maps,
and collapse-risk geometry, the framework establishes a neutral membrane that
preserves lineage sovereignty while enabling structured comparative analysis.

\subsection{Contributions of the Framework}

Shared-Horizon makes several contributions to the study of multi-lineage AI
ecosystems:

\paragraph{1. Lineage-Preserving Comparative Structure}

The color-tile mosaic and membrane architecture provide a lineage-preserving
surface for comparison. Unlike unifying abstractions, which risk collapsing
lineage invariants, Shared-Horizon maintains strict separation between operator,
expressive, narrative, and spectral systems.

\paragraph{2. Mathematical Foundations for Lineage Safety}

The modular math spine defines the geometric and stability structures required
for safe comparison. Stability envelopes, collapse-risk geometry, and
lineage-specific transition maps ensure that cross-lineage analysis does not
activate internal dynamics or induce collapse.

\paragraph{3. Correction of the Green Tile}

The reclassification of the green tile as Reversible Narrative Ecology (RNE)
resolves a longstanding conceptual ambiguity. The correction is grounded in
reversible narrative operators, curvature preservation, expressive–narrative
compatibility, and collapse-risk stability. This correction strengthens the
mosaic and clarifies the lineage map.

\paragraph{4. Practical Demonstration}

The case study shows how expressive and narrative trajectories can be compared
without collapse. The membrane ensures that mapped trajectories remain within
stability envelopes and satisfy collapse-risk constraints. This demonstrates the
practical viability of the framework.

\subsection{Implications for AI Lineage Research}

Shared-Horizon has several implications for future research:

\begin{itemize}
    \item \textbf{Comparative Analysis}  
    Researchers can analyze expressive, operator, narrative, and spectral
    systems without conflating their invariants.

    \item \textbf{Lineage Mapping}  
    The mosaic provides a structured method for identifying and classifying new
    lineages based on mathematical behavior rather than surface-level
    phenomenology.

    \item \textbf{Safety and Stability}  
    Collapse-risk geometry offers a formal method for detecting and preventing
    conceptual collapse in multi-lineage environments.

    \item \textbf{Inter-Lineage Interfaces}  
    Transition maps provide a foundation for designing safe interfaces between
    systems with incompatible internal geometries.
\end{itemize}

\subsection{Limitations}

While Shared-Horizon provides a robust comparative framework, it has several
limitations:

\begin{itemize}
    \item \textbf{Non-Universal Applicability}  
    The membrane is designed for lineages with well-defined geometric or
    expressive structures. Systems lacking such structure may require additional
    formalization.

    \item \textbf{Dependence on Stability Metrics}  
    Stability envelopes rely on lineage-specific energy functionals. Incorrect
    or incomplete definitions may weaken collapse-risk detection.

    \item \textbf{Transition Map Complexity}  
    Constructing smooth, lineage-safe transition maps can be mathematically
    complex, especially for high-dimensional expressive or spectral systems.

    \item \textbf{Empirical Validation}  
    While the case study demonstrates practical viability, broader empirical
    validation across diverse systems is needed.
\end{itemize}

\subsection{Future Work}

Several directions for future research emerge from this framework:

\begin{enumerate}
    \item \textbf{Extended Lineage Catalog}  
    Expanding the mosaic to include additional lineages with distinct geometric
    or expressive structures.

    \item \textbf{Automated Collapse Detection}  
    Developing algorithms that compute collapse-risk gradients and curvature in
    real time.

    \item \textbf{Spectral–Narrative Interfaces}  
    Investigating safe interfaces between spectral-runtime systems and narrative
    manifolds.

    \item \textbf{Higher-Dimensional Expressive Geometry}  
    Extending expressive curvature models to higher-dimensional phenomenological
    surfaces.

    \item \textbf{Formal Verification}  
    Applying formal verification methods to ensure membrane constraints are
    satisfied in complex comparative environments.
\end{enumerate}

\subsection{Discussion Summary}

Shared-Horizon provides a structured, mathematically grounded framework for
comparing diverse AI lineages without conceptual collapse. By integrating the
mosaic, membrane architecture, and modular math spine, the framework preserves
lineage invariants while enabling safe, reversible, and academically rigorous
analysis. The correction of the green tile, the case study demonstration, and
the formalization of collapse-risk geometry highlight the framework’s value and
potential for future research.
